\chapter{Tendonitis}
\index{Tendonitis}

\medskip
\noindent\textit{Educational reference; always seek professional advice for individual care.}

\section*{Definition and Overview}
Tendonitis is the traditional name for inflammation of a tendon, the strong, fibrous cord that connects a muscle to a bone. The word comes from the Latin for tendon plus the suffix "-itis", meaning inflammation, and it is the term most people use for a tendon that has become painful, swollen, and tender after an injury or a burst of overuse. Common examples include Achilles tendonitis at the back of the ankle, patellar tendonitis just below the kneecap, rotator cuff tendonitis in the shoulder, and the elbow and wrist tendon problems that follow repetitive work or sport. In recent decades, medical understanding has refined the picture: classic inflammation with immune cells is present in some acute cases, but many "tendonitis" problems are actually a mixture of mild inflammation and a painful overload reaction of the tendon tissue itself, and chronic problems are better described by the terms tendinopathy and tendinosis. Even so, the label tendonitis remains useful for the acute, recent-onset pattern of tendon pain, and it points towards sensible first aid: relative rest, ice, simple pain relief, and gentle rehabilitation. Most episodes of acute tendonitis settle within days to a few weeks with good management, but without it the problem can become chronic and much harder to treat. This chapter explains how tendonitis develops, how to recognise it and tell it apart from more serious injuries, and how to treat and prevent it.

\section*{Epidemiology}
Tendonitis is one of the commonest acute soft-tissue injuries, affecting people of all ages and both sexes, though it is particularly familiar to recreational and competitive athletes, manual workers, and people who suddenly increase their level of activity. It is most frequent in the prime activity years between the late twenties and fifties, and certain tendons dominate: the Achilles and patellar tendons in runners and jumpers, the rotator cuff tendons in swimmers, throwers, and those doing overhead work, the wrist and hand flexors in people who grip, type, or use tools, and the outer elbow tendon in racquet sports and repetitive gripping. A common pattern is the "weekend warrior" who attempts a jump in training or a burst of exercise that the tendon is not prepared for. Many episodes never reach a doctor, because they settle with rest, so the true rate is higher than clinic figures suggest. Because tendonitis is an overuse problem, its frequency rises with sports participation and with physical jobs, and repeated episodes are common when people return to full activity too quickly.

\section*{Aetiology and Pathophysiology}
Tendons transmit the pull of muscle to bone, and they are built for strength and resilience but are not invulnerable. Acute tendonitis is essentially a reaction to sudden or excessive load. When a tendon is stretched, twisted, or loaded beyond its current capacity, micro-tears develop in the collagen fibres, small blood vessels leak, and the surrounding tissues swell.

\begin{itemize}
    \item \textbf{Sudden overloading:} a jump, sprint, throw, lift, or fall that exceeds what the tendon is accustomed to, causing micro-tears
    \item \textbf{Rapid increase in training:} increasing distance, weight, intensity, or frequency faster than the tendon can adapt
    \item \textbf{Repetitive activity:} thousands of small loads in work or sport, such as running, hammering, typing, or sweeping, that accumulate faster than repair
    \item \textbf{Direct injury:} a blow or awkward twist that damages the tendon directly
    \item \textbf{Poor preparation:} inadequate warm-up, fatigue, or cold muscles training a tendon that is not ready
    \item \textbf{Biomechanical factors:} muscle weakness, tightness, poor foot posture, and faulty technique increase the load on particular tendons
    \item \textbf{Weakness and deconditioning:} after time off, a tendon is weaker than the muscle driving it and is easily overloaded
\end{itemize}

In the acute phase, the tendon and its surrounding sheath show the features of inflammation: swelling, heat, redness, and pain, driven by inflammatory chemicals and fluid. If the inflammation settles and the tendon is given appropriate graded loading, healing proceeds normally. If the tendon is repeatedly irritated, or if the person pushes through pain, the process can shift into the chronic, degenerative state that clinicians call tendinosis.

\section*{Clinical Features}
\begin{itemize}
    \item \textbf{Pain over the tendon:} pain located precisely along the tendon, such as the back of the heel, below the kneecap, over the shoulder, or at the elbow or wrist
    \item \textbf{Swelling and tenderness:} the tendon area feels swollen and is tender to touch, and the skin over it may feel warm
    \item \textbf{Pain with movement:} moving the affected muscle or stretching the tendon reproduces the pain, especially against resistance
    \item \textbf{Recent onset:} the pain typically starts after a specific event or a recent increase in activity, over days rather than months
    \item \textbf{Pain after rest:} stiffness and soreness when first moving after sitting or after waking
    \item \textbf{Reduced function:} difficulty with the aggravating activity, such as walking uphill, climbing stairs, lifting, gripping, or reaching overhead
    \item \textbf{Mild redness or heat:} in some cases the skin over the tendon is slightly red or warm, reflecting the inflammatory process
    \item \textbf{Noise with movement:} occasional creaking or grating felt over the tendon as it moves, particularly around the wrist and shoulder
\end{itemize}

The key feature distinguishing acute tendonitis from its chronic relatives is the timing: pain that is recent, relates to a clear cause, and is accompanied by swelling and tenderness points strongly to the acute inflammatory pattern.

\section*{Differential Diagnosis}
\begin{itemize}
    \item \textbf{Tendon rupture or tear:} a sudden "pop" with bruising, weakness, and a gap in the tendon indicates a tear, which needs urgent assessment; the Achilles is the classic example
    \item \textbf{Tendinosis and chronic tendinopathy:} pain that has been present for months without swelling or a clear trigger is more likely the chronic, degenerative form
    \item \textbf{Bursitis:} inflammation of the small fluid sacs near joints causes a swollen, boggy lump over a joint or bony point, and often coexists with tendonitis
    \item \textbf{Arthritis:} joint pain, stiffness, and swelling from arthritis can mimic tendon pain, particularly in the shoulder, elbow, and knee
    \item \textbf{Fracture or stress fracture:} a stress fracture or avulsion fracture causes bone tenderness and pain with weight-bearing that can be confused with tendon pain
    \item \textbf{Nerve entrapment:} nerve compression, such as carpal tunnel or ulnar nerve entrapment, causes tingling, numbness, and referred pain that can resemble tendonitis
    \item \textbf{Infection:} a rare but serious infective tendon or joint problem causes severe pain, redness, swelling, and fever and requires urgent treatment
\end{itemize}

\section*{When to Seek Medical Care}
Most acute tendonitis can be managed at home with first aid and sensible rest, but a professional opinion is wise in several situations.

Seek medical advice if you have:
\begin{itemize}
    \item Pain that does not improve within a week or two of rest and self-care
    \item Significant swelling, bruising, or weakness, which may indicate a tear
    \item Difficulty using the affected limb for daily activities
    \item Pain that wakes you at night or is present even at rest
    \item Recurrent episodes of tendon pain
\end{itemize}

\textbf{Contact your local emergency services or seek urgent care for:}
\begin{itemize}
    \item A sudden pop in a tendon followed by weakness, such as inability to stand on tiptoe after a pop in the calf or heel (Achilles rupture)
    \item Severe pain with rapidly spreading redness, swelling, heat, and fever, which can indicate infection
    \item Inability to move a limb or bear weight after an injury
\end{itemize}

\section*{Diagnosis and Investigations}
The diagnosis of tendonitis is usually straightforward and clinical, and imaging is reserved for when a tear, rupture, or another condition is suspected.

\begin{itemize}
    \item \textbf{History:} the onset and timing of pain, the activities that preceded it, recent changes in training or work, and any past tendon problems
    \item \textbf{Examination:} locating tenderness precisely over the tendon, checking for swelling and warmth, and testing the tendon with resisted movements and stretching
    \item \textbf{Functional tests:} specific movements that reproduce the pain, such as resisted grip for the elbow or a heel raise for the Achilles
    \item \textbf{Ultrasound:} shows swelling, fluid, and any tears in the tendon, and is useful when a rupture is suspected or the diagnosis is unclear
    \item \textbf{MRI:} used for deep tendons, for suspected partial tears, or when symptoms do not settle as expected
    \item \textbf{X-ray:} occasionally used to exclude a fracture or to see calcification in the tendon
    \item \textbf{Blood tests:} only when infection, gout, or inflammatory arthritis needs to be excluded
\end{itemize}

\section*{Management}
Acute tendonitis is treated first with the familiar principles of soft-tissue first aid, followed by a gentle return to activity.

\begin{itemize}
    \item \textbf{Relative rest:} stop or greatly reduce the aggravating activity for a few days, but avoid complete rest, which weakens the tendon further
    \item \textbf{Ice:} apply a cold pack to the tender area for fifteen to twenty minutes several times a day for the first couple of days to reduce swelling and pain
    \item \textbf{Compression and elevation:} a firm bandage and keeping the limb raised help limit swelling, particularly for ankle and wrist tendons
    \item \textbf{Pain relief:} paracetamol, or a short course of an anti-inflammatory medicine such as ibuprofen, helps control pain and inflammation; use them as directed and briefly
    \item \textbf{Gentle movement and stretching:} after the first few days, begin gentle range-of-movement and stretching exercises to prevent stiffness
    \item \textbf{Progressive strengthening:} once pain settles, gradually rebuild tendon strength with a graded exercise program, ideally under physiotherapy guidance
    \item \textbf{Physiotherapy:} a physiotherapist can guide rehabilitation, correct technique and biomechanics, and advise on a safe return to sport or work
    \item \textbf{Strapping and supports:} temporary taping, braces, or a heel lift can off-load the tendon during recovery
    \item \textbf{Avoid further irritation:} do not push through pain; if an activity is aggravating the tendon, stop it and modify the load
    \item \textbf{Injections:} reserved for stubborn cases and used cautiously, because corticosteroid injections weaken the tendon in the longer term
\end{itemize}

\section*{Complications}
\begin{itemize}
    \item \textbf{Progression to chronic tendinosis:} repeated irritation or a quick return to full activity converts acute tendonitis into a chronic, harder-to-treat condition
    \item \textbf{Tendon rupture:} a weakened, repeatedly injured tendon can tear, occasionally completely, and is a serious setback
    \item \textbf{Recurrence:} without proper rehabilitation, the tendon remains vulnerable and the problem tends to return
    \item \textbf{Secondary problems:} favouring the painful tendon can strain other muscles and joints, causing new pain
    \item \textbf{Long-term stiffness and weakness:} inadequate rehabilitation leaves the tendon shortened, weak, and prone to re-injury
\end{itemize}

\section*{Prognosis and Outcomes}
The outlook for genuine acute tendonitis is very good. With relative rest, ice, simple pain relief, and a sensible return to activity, most people are noticeably better within a few days to two weeks and fully recovered within a few weeks to a couple of months. The two factors that most often derail recovery are returning to full activity too soon and doing nothing at all: the first re-injures the tendon, and the second lets it stiffen and weaken. A graduated return, guided by pain rather than by the calendar, gives the best results. Where an episode has been severe or recurrent, the tendon may take longer, and a course of physiotherapy to rebuild strength and correct the underlying cause of the overload is the most reliable route back. The key message is that tendonitis is a common, treatable injury that responds well to commonsense first aid and a patient return to activity.

\section*{Prevention and Self-Care}
\begin{itemize}
    \item Warm up properly and build up the intensity and duration of activity gradually, especially after time off
    \item Strengthen the muscles that move commonly affected tendons, and stretch the tight muscles around them
    \item Use correct technique and well-fitting footwear, and address poor foot posture and muscle imbalance with professional help
    \item Take rest days and avoid a sudden jump in training volume or weight
    \item Modify repetitive work tasks, with regular breaks and ergonomic changes where needed
    \item At the first sign of tendon pain, reduce the aggravating activity and start simple first aid rather than pushing through
    \item Return to full activity gradually, guided by pain, after any tendon injury
\end{itemize}

\section*{Sources and Further Reading}
The following organisations provide reliable information about tendonitis and acute tendon injuries.

\begin{itemize}
    \item Sports Medicine Australia. \textit{Tendon injuries: first aid and return to sport.} https://sma.org.au/
    \item healthdirect Australia. \textit{Tendonitis: symptoms and treatment.} https://www.healthdirect.gov.au/
    \item NHS. \textit{Tendinopathy and tendon injuries.} https://www.nhs.uk/
    \item Mayo Clinic. \textit{Tendinitis: symptoms, causes, and treatment.} https://www.mayoclinic.org/
    \item Australian Physiotherapy Association. \textit{Acute tendon injuries: a recovery guide.} https://australian.physio/
    \item Better Health Channel (Victoria). \textit{Tendon injuries.} https://www.betterhealth.vic.gov.au/
\end{itemize}

\clearpage
